% Compile with: latexmk -cd -pdf -interaction=nonstopmode -outdir=. acceptance_letter.tex
\documentclass[conference]{IEEEtran}

\usepackage{amsmath, amssymb} % Useful for equations
\usepackage{graphicx} % For including images
\usepackage{cite} % For bibliography management
\usepackage{array} % For p{} column width control
\usepackage{xcolor}
\usepackage{tcolorbox}

\date{\today}

\begin{document}


\title{Towards a Practical Framework for Cryptographic Protocols in Adversarial Cooperation}
\author{
    \IEEEauthorblockN{S. Restrepo Ruiz}
    \IEEEauthorblockA{waajacu.com \\ Email: santiago.restrepo.ruiz@gmail.com}
}

\maketitle

\begin{abstract}
This study aims to investigates the role of information theory in designing cryptographic protocols in adversarial settings. 
By leveraging secure multi-party computations; during the research we will build 
an \textit{open-source} software library demonstrating the results.
\end{abstract}

\section{Introduction}
To de-escalate conflict, we propose to have adversaries cooperate in a structured setting. Where trust is cryptographically enforced. 
---Implement information-secure and verifiable encryption protocols, securing game-theoretic coordination, and privacy-preserving shared decision-making. 

\section{Literature Review}
Cryptography extends far beyond Encryption, it serves as the foundation of digital trust. The tools at hand for this study are: 
\textcolor{blue}{\textit{Zero-knowledge proofs}}, 
\textcolor{blue}{\textit{Homomorphic Computation}}, 
\textcolor{blue}{\textit{Garbled Circuits}}, 
\textcolor{blue}{\textit{Oblivious Transfer}}, 
\textcolor{blue}{\textit{Blockchains}}, 
\textcolor{blue}{\textit{Differential Privacy}}, and 
\textcolor{blue}{\textit{Threshold Cryptography}}; 
---Combining these secure units enable broad collaborative problem-solving that are believed to be key in adversarial cooperation.

\vspace{5pt}
\section{Research Questions}

\begin{tabular}{lp{210pt}}
    {\small \textbf{1)}} & {\small What does it mean to apply cryptography to cooperation?} \\
    {\small \textbf{2)}} & {\small How can the proposed cryptographic protocols be improved through interdisciplinary collaboration?}
\end{tabular}
    
\vspace{5pt}
\section{Research Objectives}

\begin{tabular}{lp{210pt}}
    {\small \textbf{1)}} & {\small Demonstrate for classical game-theoretic problems.} \\
    {\small \textbf{2)}} & {\small Open-source a software library based on this research.} 
\end{tabular}

\section{Methodology}
\vspace{5pt}
The study will focus on the following structured games:
\vspace{5pt}

\begin{tabular}{lp{210pt}}
    {\small \textbf{1)}} & {\textcolor{blue}{\textit{Playing card games without revealing one's cards:}} \linebreak---Demonstrate how to claim victory in a card game without exposing the cards at hand. } \\
    {\small \textbf{2)}} & {\textcolor{blue}{\textit{Playing tic-tac-toe without revealing one's strategy:}} \linebreak---Demonstrate how a player can prove they have a winning strategy in tic-tac-toe without disclosing the actual moves. } \\
    {\small \textbf{3)}} & {\textcolor{blue}{\textit{Classical adversarial stag-hunt:}} ---Demonstrate a path to equilibrium for this fundamental game of adversarial cooperation. } 
\end{tabular}

\section{Acceptance Letter Request and Justification}
I kindly, respectfully ask the distinguished faculty members consideration and support for my application to \textit{Beihang University} Master's in Cyberspace Security. 

Universities in Colombia currently lack research programs in Applied Cryptography. I have been diligently studying the Chinese language for several years, 
recognizing the importance of seamless communication and collaboration in a Chinese academic setting. 
I am deeply enthusiastic about the prospect of learning from and contributing to the vibrant research community in China. 
I believe this scholarship would not only enable me to pursue my professional aspirations but also foster valuable academic exchange between our nations.

% \section{Target Universities and Programs}
% \vspace{5pt}
% {\small English-taught program:} \\

% \begin{tabular}{rl}
%     \textcolor{blue}{\small\textit{Beihang University}} & {\small Master's Cyberspace Security}\\
%     % \textcolor{blue}{\small\textit{Beihang University}} & {\small Master's Integrated Circuit Engineering}
% \end{tabular}

\section{Expected Outcomes}

\begin{tabular}{lp{210pt}}
    {\small \textbf{1)}} & {\small Demonstration adversarial cooperation in game-theory.} \\
    {\small \textbf{2)}} & {\small Development of an open-source cryptographic protocol suite enabling cooperation in adversarial settings.} \\
    {\small \textbf{3)}} & {\small Advocate on how cryptography facilitate trust in conflict.}
\end{tabular}

\section{Conclusion}
This research seeks to show how cryptographic techniques can be used for peace-building and trust establishment. 

\begin{thebibliography}{99}

\bibitem{Yao1982}
A. C. Yao, ``Protocols for secure computations,'' in *Proceedings of the 23rd Annual Symposium on Foundations of Computer Science*, 1982.

\bibitem{Goldwasser1989}
S. Goldwasser, S. Micali, and C. Rackoff, ``The Knowledge Complexity of Interactive Proof Systems,'' in *SIAM Journal on Computing*, 1989.

\bibitem{Gentry2009}
C. Gentry, ``Fully Homomorphic Encryption Using Ideal Lattices,'' in *Proceedings of the 41st ACM Symposium on Theory of Computing (STOC)*, 2009.

\bibitem{Dwork2006}
C. Dwork, ``Differential Privacy,'' in *Proceedings of the International Colloquium on Automata, Languages, and Programming (ICALP)*, 2006.

\bibitem{Canetti2000}
R. Canetti, ``Security and Composition of Multi-Party Cryptographic Protocols,'' in *Journal of Cryptology*, 2000.

\bibitem{Rosen2004}
A. Rosen, ``Concurrent Zero-Knowledge,'' in *Springer*, 2004.

\bibitem{Menezes1996}
A. J. Menezes, P. C. van Oorschot, and S. A. Vanstone, *Handbook of Applied Cryptography*, CRC Press, 1996.

\bibitem{Zhu2021}
T. Zhu, G. Li, W. Zhou, and P. Xiong, *Differential Privacy and Applications*, Springer, 2021.

\bibitem{Cramer2015}
R. Cramer, I. Damgård, and J. B. Nielsen, *Secure Multiparty Computation and Secret Sharing: An Information-Theoretic Approach*, Cambridge University Press, 2015.

\end{thebibliography}

\bibliographystyle{IEEEtran}
\bibliography{references}

\end{document}
